\section*{Capítulo \ref{ch:g2:numbers} --- Números até 1000}

\begin{solution}{exo:g2:numbers:1}
$258$: $2$ centenas, $5$ dezenas, $8$ unidades. $470$: $4$, $7$, $0$.
$600$: $6$, $0$, $0$. $909$: $9$, $0$, $9$.
\end{solution}

\begin{solution}{exo:g2:numbers:2}
$407$; \ $630$; \ $1000$.
\end{solution}

\begin{solution}{exo:g2:numbers:3}
$216$; \ $504$; \ $790$.
\end{solution}

\begin{solution}{exo:g2:numbers:4}
$438 = 400 + 30 + 8$; \quad $702 = 700 + 2$; \quad
$560 = 500 + 60$.
\end{solution}

\begin{solution}{exo:g2:numbers:5}
$327 < 372$; \ $580 > 508$; \ $199 < 200$; \ $640 > 604$.
\end{solution}

\begin{solution}{exo:g2:numbers:6}
$304 < 340 < 403 < 430$.
\end{solution}

\begin{solution}{exo:g2:numbers:7}
$253 + 10 = 263$; \quad $253 + 100 = 353$; \quad
$480 - 10 = 470$; \quad $716 - 100 = 616$.
\end{solution}

\begin{solution}{exo:g2:numbers:8}
$195, 205, 215, 225, 235$ (o primeiro salto preenche as dezenas e
carrega). $300, 400, 500, 600, 700$.
\end{solution}

\begin{solution}{exo:g2:numbers:9}
$6$ centenas e $5$ unidades: $600 + 5 = 605$ lápis.
\end{solution}

\begin{solution}{exo:g2:numbers:10}
O maior: $852$. Menor: $258$. Distância:
$852 - 258 = 594$.
\end{solution}
